\chapter{Sistema anti-bunny-virus}

\section{Arquitectura y funcionamiento general}

\emph{Anti-bunny-virus} es un monitor de espacio de usuario. Cada 0,25 s obtiene
procesos y recursos desde \texttt{/proc}, mide un directorio, evalúa reglas y
registra o contiene la amenaza. Al iniciar configura el log, el modo de respuesta
y la autorización del laboratorio. Usa arreglos estáticos para un máximo de
8192 procesos y recursos, 1024 archivos y 256 alertas por ciclo.

\paragraph{Organización del código.}
El repositorio separa recolectores (\texttt{monitor\_procesos},
\texttt{monitor\_recursos}, \texttt{monitor\_archivos}), el motor
(\texttt{motor\_deteccion}), la respuesta (\texttt{respuesta}) y el ciclo
principal (\texttt{main.c}). La versión integrada desde \texttt{main} mantiene
el mismo comportamiento evaluable, pero documenta en comentarios el propósito de
cada función pública y de los pasos críticos (historial por
\texttt{starttime}, persistencia, atribución y contención). Los simuladores en
\texttt{simuladores/} incluyen la misma documentación breve sobre la carga que
generan.

\begin{figure}[htbp]
  \centering
  \fbox{\parbox{0.78\textwidth}{\centering
    Recolección de procesos, CPU/RSS y archivos}}
  \par\smallskip$\Downarrow$\par\smallskip
  \fbox{\parbox{0.78\textwidth}{\centering
    Historial por PID--\texttt{starttime} y por ruta}}
  \par\smallskip$\Downarrow$\par\smallskip
  \fbox{\parbox{0.78\textwidth}{\centering
    Tasas temporales, umbrales y persistencia}}
  \par\smallskip$\Downarrow$\par\smallskip
  \fbox{\parbox{0.78\textwidth}{\centering
    Alerta con severidad, identidad y evidencia}}
  \par\smallskip$\Downarrow$\par\smallskip
  \fbox{\parbox{0.78\textwidth}{\centering
    Registro o contención del grupo autorizado}}
  \caption{Flujo principal del sistema anti-bunny-virus.}
  \label{fig:flujo-sistema}
\end{figure}

El Listado~\ref{alg:ciclo-principal} resume \texttt{main.c}. Todas las alertas
se registran; las señales solo se envían bajo la política de la
Sección~\ref{sec:deteccion-respuesta}.

\Needspace{18\baselineskip}
\begin{lstlisting}[style=pseudocodigo,caption={Funcionamiento general de anti-bunny-virus.},label={alg:ciclo-principal}]
configurar log, modo y autorizacion de laboratorio
instalar manejador de SIGINT
inicializar historiales de procesos, recursos y archivos

mientras la ejecucion siga activa:
    procesos <- recolectar procesos desde /proc
    recursos <- recolectar CPU y RSS desde /proc
    archivos <- medir archivos del directorio vigilado

    comparar cada muestra con su historial
    calcular hijos nuevos y CPU
    calcular crecimiento_RAM en MB/s
    calcular crecimiento_archivo en MB/s

    si existen archivos:
        mayor <- archivo con mayor crecimiento_archivo
        si crecimiento_archivo(mayor) > 0:
            atribuir mayor a un proceso, si es posible

    alertas <- detectar(procesos, recursos, archivos)
    registrar y responder(alertas, modo)
    actualizar historiales
    esperar 0.25 segundos
\end{lstlisting}

\section{Monitorización y extracción de métricas}

\paragraph{Procesos.}

\texttt{monitor\_procesos} recorre los PID de \texttt{/proc} y obtiene los datos
de la Tabla~\ref{tab:campos-proceso}. Si un proceso termina durante la lectura,
descarta esa muestra.

\begin{table}[H]
  \centering
  \small
  \caption{Datos recolectados para cada proceso.}
  \label{tab:campos-proceso}
  \begin{tabular}{p{0.22\textwidth}p{0.31\textwidth}p{0.36\textwidth}}
    \toprule
    \textbf{Dato} & \textbf{Fuente} & \textbf{Uso} \\
    \midrule
    PID, PPID y PGID & \texttt{/proc/[pid]/stat} & Identidad, relación padre--hijo y grupo de contención. \\
    Tiempo de inicio & \texttt{/proc/[pid]/stat} & Diferencia procesos que reutilizan el mismo PID. \\
    UID y usuario & Metadatos de \texttt{/proc/[pid]} & Verifica el propietario autorizado. \\
    Nombre y ejecutable & \texttt{stat} y \texttt{exe} & Describe la evidencia de la alerta. \\
    Hijos & Comparación de PPID & Detecta proliferación por proceso padre. \\
    \bottomrule
  \end{tabular}
\end{table}

El módulo cuenta los hijos por PPID. El historial identifica cada proceso por
PID y tiempo de inicio para evitar confusiones si el PID se reutiliza. Los hijos
nuevos y su tasa son

\begin{equation}
  H_{\mathrm{nuevos}}(p,t)=\max(0,H(p,t)-H(p,t-1)),
\end{equation}

\begin{equation}
  R_H(p,t)=\frac{H_{\mathrm{nuevos}}(p,t)}{\Delta t},
\end{equation}

donde $\Delta t$ proviene de un reloj monotónico. El historial se actualiza al
final de cada ciclo.

\paragraph{CPU y memoria.}

\texttt{monitor\_recursos} lee los ticks de CPU de
\texttt{/proc/[pid]/stat} y la memoria residente \texttt{VmRSS} de
\texttt{/proc/[pid]/status}. También identifica el historial mediante PID y
tiempo de inicio.

Sea $C_t$ la suma de ticks de CPU en la muestra actual, $f$ la frecuencia de
ticks del sistema y $M_t$ la memoria residente en MB. Las métricas implementadas
son

\begin{equation}
  \mathrm{CPU}(t)=\frac{(C_t-C_{t-1})/f}{\Delta t}\times 100,
\end{equation}

\begin{equation}
  \Delta\mathrm{RSS}(t)=\frac{M_t-M_{t-1}}{\Delta t}.
\end{equation}

La primera muestra produce tasas iguales a cero. La CPU calculada corresponde
al proceso durante el intervalo, no a la carga global de la máquina virtual.

\paragraph{Archivos y atribución.}

El monitor recorre \path{/tmp/anti_bunny_demo} y conserva el tamaño de cada
ruta. Si $S_t$ es el tamaño actual e $I$ el intervalo, calcula

\begin{equation}
  R_A(t)=\frac{\max(0,S_t-S_{t-1})}{1024^2 I}\quad\mathrm{MB/s}.
\end{equation}

Una disminución o una ruta nueva produce tasa cero. En cada ciclo solo se intenta
atribuir el archivo de mayor crecimiento.

La atribución compara la ruta con los enlaces de \texttt{/proc/[pid]/fd}, como
muestra el Listado~\ref{lst:atribucion}. Conserva el primer PID y PGID
coincidente. Sin coincidencia puede existir alerta, pero no contención.

\begin{lstlisting}[style=codigoC,caption={Comparación usada para atribuir un archivo abierto.},label={lst:atribucion}]
longitud = readlink(ruta_fd, enlace, sizeof(enlace) - 1);
if (longitud < 0) continue;
enlace[longitud] = '\0';
if (strcmp(enlace, archivo->ruta) == 0) {
    archivo->pid_escritor = procesos[i].pid;
    archivo->pgid_escritor = procesos[i].pgid;
    archivo->atribuido = 1;
    closedir(fds);
    return;
}
\end{lstlisting}

La atribución solo identifica archivos abiertos durante el sondeo; no constituye
evidencia forense de una escritura concreta.

\section{Detección y respuesta}
\label{sec:deteccion-respuesta}

\paragraph{Motor de detección.}

El motor identifica cada estado por tipo, PID y tiempo de inicio. Cuenta ciclos
anómalos y solo emite otra alerta si aumenta la severidad; al normalizarse,
reinicia el estado. La Tabla~\ref{tab:umbrales} contiene los umbrales.

\begin{table}[H]
  \centering
  \small
  \caption{Parámetros de detección configurados.}
  \label{tab:umbrales}
  \begin{tabular}{p{0.43\textwidth}p{0.18\textwidth}p{0.24\textwidth}}
    \toprule
    \textbf{Parámetro} & \textbf{Valor} & \textbf{Aplicación} \\
    \midrule
    Intervalo de monitorización & 0,25 s & Todos los recolectores. \\
    Hijos por proceso & 20 & Regla de proliferación. \\
    Hijos totales o nuevos por ventana & 8 & Inicio temprano de proliferación. \\
    Persistencia crítica & 2 ciclos & Procesos, recursos y archivos. \\
    Memoria residente & 256 MB & Regla de recursos. \\
    Crecimiento de RSS & 32 MB/s & Regla de recursos y correlación. \\
    CPU del proceso & 80\% & Regla de recursos y correlación. \\
    Crecimiento de archivo & 5 MB/s & Regla de archivos. \\
    Gracia de terminación & 2 s & Escalamiento de señales. \\
    \bottomrule
  \end{tabular}
\end{table}

Los dos umbrales expresados en MB/s miden objetos distintos:
\emph{crecimiento de RSS} es el aumento de RAM ocupada por un proceso, mientras
que \emph{crecimiento de archivo} es el aumento de una ruta en disco.

La proliferación genera una alerta sospechosa y pasa a crítica y accionable tras
dos ciclos o al alcanzar 40 hijos. CPU, RSS y archivos solo generan alertas no
accionables. Una alerta de archivo puede ser crítica si su escritor también
supera CPU o crecimiento de RSS. Por tanto, solo la proliferación puede activar
contención.

El Listado~\ref{lst:regla-procesos} muestra la regla de proliferación en C.

\Needspace{13\baselineskip}
\begin{lstlisting}[style=codigoC,caption={Clasificación de la proliferación de procesos.},label={lst:regla-procesos}]
int activo = procesos[i].hijos >= max_hijos ||
             procesos[i].hijos >= max_hijos_nuevos ||
             procesos[i].hijos_nuevos >= max_hijos_nuevos;

int critica = ciclos >= ciclos_persistencia ||
              procesos[i].hijos >= max_hijos * 2;

agregar(alertas_out, &total, max_alertas,
        procesos[i].pid, procesos[i].pgid,
        procesos[i].starttime, procesos[i].uid,
        critica, critica ? "critica" : "sospechosa",
        "procesos",
        critica ? "Proliferacion de procesos persistente"
                : "Proliferacion de procesos observada",
        evidencia);
\end{lstlisting}

\Needspace{32\baselineskip}
El Listado~\ref{alg:deteccion} resume las tres reglas con nombres explícitos
para RAM y archivos.

\begin{lstlisting}[style=pseudocodigo,caption={Motor de detección y clasificación.},label={alg:deteccion}]
para cada proceso p:
    si hijos_totales(p) >= 8 o hijos_nuevos(p) >= 8:
        aumentar ciclos_anomalos(p)
        si ciclos_anomalos(p) >= 2 o hijos_totales(p) >= 40:
            emitir alerta CRITICA y ACCIONABLE
        si no:
            emitir alerta SOSPECHOSA y NO ACCIONABLE
    si no:
        reiniciar ciclos_anomalos(p)

para cada recurso r:
    RAM_actual <- RSS(r) en MB
    crecimiento_RAM <- cambio de RSS por segundo
    si CPU(r) > 80 o RAM_actual > 256
       o crecimiento_RAM > 32 MB/s:
        aumentar ciclos_anomalos(r)
        si ciclos_anomalos(r) >= 2:
            emitir alerta SOSPECHOSA y NO ACCIONABLE
    si no:
        reiniciar ciclos_anomalos(r)

para cada archivo a:
    crecimiento_archivo <- cambio de tamano por segundo
    si crecimiento_archivo > 5 MB/s:
        aumentar ciclos_anomalos(a)
        si ciclos_anomalos(a) >= 2:
            si el escritor supera CPU o crecimiento_RAM:
                emitir alerta CRITICA y NO ACCIONABLE
            si no:
                emitir alerta SOSPECHOSA y NO ACCIONABLE
    si no:
        reiniciar ciclos_anomalos(a)
\end{lstlisting}

\paragraph{Registro y contención.}

Cada alerta se escribe en la salida estándar y en \texttt{logs/eventos.log} con
una línea \texttt{WARNING} que incluye fecha, severidad, tipo, PID, PGID,
\texttt{starttime}, descripción, evidencia y modo. Las acciones de contención o
denegación se registran como líneas \texttt{INFO} o \texttt{ERROR} con el mismo
formato de campos clave--valor.

El modo predeterminado \texttt{alerta} solo registra. En
\texttt{terminar\_lab}, la alerta debe ser crítica y accionable; PID y PGID deben
ser mayores que uno; y PGID, UID y tiempo de inicio deben coincidir con
\texttt{/proc} y la autorización. El Listado~\ref{lst:identidad} muestra la
revalidación.

\begin{lstlisting}[style=codigoC,caption={Revalidación de la identidad antes de contener (\texttt{respuesta.c}).},label={lst:identidad}]
if (stat(ruta, &datos) != 0) return 0;
pgid_actual = getpgid(a->pid);

return starttime_actual == a->starttime &&
       pgid_actual == a->pgid &&
       datos.st_uid == uid_laboratorio;
\end{lstlisting}

Tras autorizar, envía \texttt{SIGTERM} al PGID y espera dos segundos. Si el grupo
sigue activo y supera una nueva validación, envía \texttt{SIGKILL}. Toda acción
o denegación queda registrada.

El Listado~\ref{alg:contencion} resume esta secuencia.

\Needspace{16\baselineskip}
\begin{lstlisting}[style=pseudocodigo,caption={Contención gradual del grupo autorizado.},label={alg:contencion}]
para cada alerta:
    escribir alerta y evidencia en el log

    si modo != terminar_lab:
        continuar
    si alerta no es critica o no es accionable:
        registrar denegacion y continuar
    si PID, PGID, UID o starttime no coinciden:
        registrar denegacion y continuar

    enviar SIGTERM al grupo autorizado
    esperar 2 segundos
    comprobar que el grupo siga activo
    revalidar PGID, UID y starttime
    si continua activo y autorizado:
        enviar SIGKILL al mismo grupo
\end{lstlisting}

\Needspace{22\baselineskip}
\section{Simuladores y ejecución controlada}

\paragraph{Configuración y ejecución.}

Los umbrales están en \texttt{src/config.h}. El modo y la ruta del log admiten
variables de entorno. El Listado~\ref{lst:ejecucion-local} muestra la ejecución
local.

\begin{lstlisting}[style=comandos,caption={Compilación y ejecución local en modo de alerta.},label={lst:ejecucion-local}]
make check

# Terminal 1: solo registra; no envia senales
ANTI_BUNNY_MODE=alerta ./anti-bunny-virus

# Terminal 2: ejecutar un solo simulador por corrida
./simuladores/simulador_fork_bomb
./simuladores/simulador_memoria
./simuladores/simulador_archivo
\end{lstlisting}

\paragraph{Simuladores incluidos.}

Debe ejecutarse un simulador por corrida y solo dentro de la máquina virtual. La
Tabla~\ref{tab:simuladores} resume la carga y el resultado esperado.

\begin{table}[H]
  \centering
  \small
  \caption{Simuladores controlados del repositorio.}
  \label{tab:simuladores}
  \begin{tabular}{p{0.17\textwidth}p{0.32\textwidth}p{0.39\textwidth}}
    \toprule
    \textbf{Simulador} & \textbf{Carga ejecutada} & \textbf{Resultado esperado} \\
    \midrule
    Procesos & Hasta 64 procesos; ocho hijos por proceso; permanencia de 5 s. & Alerta de proliferación; crítica y accionable al persistir. \\
    Memoria & Diez asignaciones de 32 MB; máximo nominal de 320 MB. & Alerta sospechosa por RSS o crecimiento de memoria. \\
    Archivo & Ocho escrituras de 2 MB cada 0,2 s; total de 16 MB. & Alerta por crecimiento; incluye ruta y escritor si se atribuye. \\
    \bottomrule
  \end{tabular}
\end{table}

Solo el simulador de procesos puede activar contención. systemd lo aísla con
\texttt{setsid}, publica su PGID y limita el servicio a 256 PID. Tras instalar
\texttt{deploy/}, se ejecuta como indica el
Listado~\ref{lst:ejecucion-protegida}.

\begin{lstlisting}[style=comandos,caption={Ejecución del escenario protegido con systemd.},label={lst:ejecucion-protegida}]
sudo systemctl start anti-bunny-virus.service
sudo systemctl start anti-bunny-attack.service
sudo journalctl -u anti-bunny-virus.service -f
\end{lstlisting}

Memoria y archivo solo validan detección; no solicitan terminación. El simulador
de archivo escribe \path{/tmp/anti_bunny_demo/archivo_crecimiento_demo.bin}.
Los tres simuladores son acotados.
